\chapter{Answers to the checkpoints}
\label{app:answers}

\small
\newcommand{\ans}[2]{\par\textbf{#1}\quad #2\par}
\setlength{\parskip}{0.45em}

\subsection*{Chapter 1}
\ans{Q1.1}{You never see the test data, and a crash anywhere means no score at all, discovered only hours later. In a
normal competition you upload a file you can check first.}
\ans{Q1.2}{Not necessarily. With \ca 390 public studies the score wobbles by about $\pm$0.005, so 0.003 is within noise.}
\ans{Q1.3}{No. +0.01 AUC is worth about 0.022 efficiency points; one hour costs about 0.11.}

\subsection*{Chapter 2}
\ans{Q2.1}{No. ``Small'' is below the threshold; only moderate or large effusions count.}
\ans{Q2.2}{Medial OA, Lateral OA and PF OA: the same cartilage-loss rule (about 1\,cm or more with over half the thickness
lost), each in a different compartment.}
\ans{Q2.3}{Medial and lateral swap sides between a left and a right knee. A side-blind model cannot use position to tell
the two menisci apart and has to rely on shape and context.}

\subsection*{Chapter 3}
\ans{Q3.1}{On T1, edema is greyish against bright fatty marrow: hard to see. On fat-suppressed PD the marrow is dark and
the edema bright: easy.}
\ans{Q3.2}{MRI brightness has no units; it depends on scanner, coil, sequence and patient. Each series must be
normalised on its own.}
\ans{Q3.3}{For example the MCL, which runs along the side of the knee and is best seen on coronal slices. (PF OA and
Baker's cyst, best on axial, are other good answers.)}

\subsection*{Chapter 4}
\ans{Q4.1}{File names are random IDs. Only the physical position gives the true order.}
\ans{Q4.2}{Not necessarily. The gold set was chosen to contain many positives, so its prevalence is higher by design.
Compare labels with gold by AUC, not by prevalence.}
\ans{Q4.3}{The test studies have no reports. Only images are available at test time.}

\subsection*{Chapter 5}
\ans{Q5.1}{Three different neighbouring slices give the network some depth context (a tear that continues across
slices) at no extra cost; a copied slice gives it nothing new.}
\ans{Q5.2}{$p = 0.5$. The model learns ``don't be confident about this one'', which is the honest message of an unknown
label.}
\ans{Q5.3}{Only a little. Same images and labels lead to strongly correlated mistakes; ensembles help when the members
differ.}

\subsection*{Chapter 6}
\ans{Q6.1}{1.0. Every sick study scores above every healthy one; only the order matters.}
\ans{Q6.2}{Nothing. The order of the predictions does not change.}
\ans{Q6.3}{Synovitis: +0.04/12 $\approx$ +0.0033 on the macro AUC, against +0.01/12 $\approx$ +0.0008 for Baker's.}

\subsection*{Chapter 7}
\ans{Q7.1}{``Absent'' means the report says the structure is normal; ``not mentioned'' is silence. For Synovitis the
difference is large: 15 of 41 silent gold studies were positive.}
\ans{Q7.2}{Near-perfect labels do not make the images easier. The image model still has to see the tear, with limited
resolution and capacity; and 58 gold studies measure it with noise.}
\ans{Q7.3}{Spending many GPU hours labelling everything with a labeller that turns out worse. The gold studies are
labelled first, and the long run starts only if they are clearly better.}

\subsection*{Chapter 8}
\ans{Q8.1}{Brightness differs between series, even within one exam; normalising per series removes those differences,
a dataset-wide rule would keep them.}
\ans{Q8.2}{$200/140 \approx 1.43$ times finer (for example 0.78\,mm $\rightarrow$ 0.55\,mm per pixel).}
\ans{Q8.3}{The study would describe an impossible knee: the sagittal series says ``left knee'', the coronal and axial
say ``right knee''. The model would see contradictory geometry.}

\subsection*{Chapter 9}
\ans{Q9.1}{8 studies $\times$ 3 series $\times$ 8 triplets = 192 images.}
\ans{Q9.2}{Gold studies are never trained on, so all 5 fold models are fair predictors and their average is a free
ensemble. Each training study was held out by exactly one fold model; only that model is unbiased for it.}
\ans{Q9.3}{Neither number alone is decisive at this noise level. Keep both and blend them, which is what we did (gold 0.903).}

\subsection*{Chapter 10}
\ans{Q10.1}{Any difference (resizing, orientation) changes what the model sees; the model was never trained on such
images, and the score drops without any error message.}
\ans{Q10.2}{The ranks cancel out: every study ends near the middle and the finding becomes close to random (AUC near 0.5
if the two models are exactly opposite).}
\ans{Q10.3}{It restores the backup of the public notebook's staged predictions, so the public result is submitted unchanged.}

\subsection*{Chapter 11}
\ans{Q11.1}{A threshold for continuing an experiment, fixed in advance, so that noise is not talked into ``progress''
after the fact.}
\ans{Q11.2}{In their model the side is used explicitly (side-aware boxes and box pooling), so a consistent orientation
helps. Our model had no use for the side, and lost the regularising effect of mirror augmentation.}
\ans{Q11.3}{For example: the public notebook's per-finding weights had been tuned on the public leaderboard and did not
carry over to the private part; or simply noise. The winners picked models with a ruler, not by probing the leaderboard.}

\subsection*{Chapter 12}
\ans{Q12.1}{AUC ignores calibration. A label set with many confidently wrong labels can rank well but teaches the model
to be confidently wrong; log loss and confident errors catch that.}
\ans{Q12.2}{Reports rarely describe synovitis, and when they do the agent would just copy the wording. Without it, the
agent has to judge from the images, like the gold radiologists.}
\ans{Q12.3}{A guide learned on other folds is applied to studies it has never seen: an honest test. The in-sample guide
has seen the answers.}

\subsection*{Chapter 13}
\ans{Q13.1}{To predict the embedding, the image features must encode what the report talks about. At test time the head
is simply removed, so no report is needed.}
\ans{Q13.2}{Mostly from the labels, plus input and backbone: their plain per-slice model already reached 0.932 on gold,
against our 0.904; then bigger inputs and ensembles on the leaderboard.}
\ans{Q13.3}{58 gold studies are too few ($\pm$0.02) and sometimes disagreed with the leaderboard. A ruler that tracks
20--30 leaderboard scores is known to measure what matters.}

\subsection*{Chapter 14}
\ans{Q14.1}{Checking the rules thread on day one: it would have allowed image-reading labels from the start, the largest
single source of the gap.}
\ans{Q14.2}{Yes, more slowly: an open-weights vision-language model running on Kaggle, or a few hundred studies annotated
by hand or by simple anatomy heuristics, followed by a detector trained on them.}
